\section{Conclusion}


Astle.js demonstrates that domain-specific state management significantly reduces development complexity for commerce applications. By encoding commerce primitives—carts, checkout flows, customer profiles—directly into the SDK, we eliminate boilerplate while preserving the benefits of reactive, unidirectional architectures.

The event-sourced foundation enables powerful capabilities: time-travel debugging, offline support, and deterministic state reconstruction. These capabilities translate directly to improved developer experience and user satisfaction.

Future work will extend Astle.js with machine learning-driven prefetching (predicting likely cart additions) and integration with emerging state management patterns (React Server Components, Suspense).

\begin{thebibliography}{9}

\bibitem{flux2014}
Facebook, ``Flux: Application Architecture for Building User Interfaces,'' \textit{Facebook Engineering Blog}, 2014.

\bibitem{redux2015}
D. Abramov, ``Redux: Predictable State Container for JavaScript Apps,'' \textit{GitHub}, 2015.

\bibitem{rxjs2015}
B. Lesh, ``RxJS: Reactive Extensions for JavaScript,'' \textit{GitHub}, 2015.

\bibitem{mobx2016}
M. Weststrate, ``MobX: Simple, Scalable State Management,'' \textit{GitHub}, 2016.

\bibitem{apollo2016}
Meteor Development Group, ``Apollo Client: A Fully-Featured GraphQL Client,'' \textit{GitHub}, 2016.

\bibitem{cqrs2010}
M. Fowler, ``CQRS,'' \textit{martinfowler.com}, 2010.

\bibitem{eventsourcing2005}
M. Fowler, ``Event Sourcing,'' \textit{martinfowler.com}, 2005.

\end{thebibliography}

